\section{Positive-dimensional actual components save one degree power}
\label{app:higher-positive-components}

Appendix~\ref{app:isolated-sharpness} proves that isolated actual
solutions can require a $D^{d+1}$ count. Positive-dimensional actual
components satisfy a smaller bound, extending Appendix~\ref{app:actual-first-order}
to every fixed derivative order. We retain the projective degree in
coefficient space throughout.

\begin{theorem}[Degrees by actual component dimension]
\label{thm:higher-actual-components}
Use the hypotheses and notation of Theorem~\ref{thm:isolated-upper},
and additionally put $T=b+\tau H$. Let $J_j$ be the sum of the
projective degrees of the $j$-dimensional irreducible components of
the reduced closure of the actual regular joint solution locus.
All components have dimension at most $d+1$, and
\[
 J_0\le qK^{d+1},\qquad
 J_j\le qK^{d+1-j}T^{j-1}\quad(1\le j\le d+1).
\]
Thus every positive-dimensional stratum has total degree
$O_{d,B,H}(D^d)$. For an ordinary equation with the challenge fixed,
let $I_j$ count the degrees of its actual regular components of
dimension $j$. Then
\[
 I_0\le B(Bb)^d,\qquad
 I_j\le B(Bb)^{d-j}b^{j-1}\quad(1\le j\le d).
\]
In particular the positive-dimensional part of a reduced generic
challenge fiber has degree $O_{d,B}(D^{d-1})$. This last assertion
excludes the isolated points of that fiber.
\end{theorem}
\begin{proof}
The initial-jet description and Taylor reconstruction are as in
Theorem~\ref{thm:isolated-upper}. They give dimension at most $d+1$
and the stated $J_0$ bound. We explain why a special family of linear
sections counts the full actual coefficient-space degree.

\paragraph{Evaluation projections count actual degrees.}
Let $C$ be an actual $j$-dimensional component, defined before any
Taylor center is adjoined. Choose independent generic parameters
$t_i,\lambda_i$ for $0\le i<j$, and use the affine linear forms
\[
 \ell_i(P,z)=P(t_i)+\lambda_i z.
\]
At a projective point at infinity, the coefficient direction of
$P$ and the $z$ direction cannot all be zero. Consequently
$P(t)+\lambda z$ is not the zero polynomial in $(t,\lambda)$ at
that point. The incidence of a point of $C$ at infinity and $j$
vanishing forms has dimension at most
$(j-1)+j=2j-1$, less than the $2j$-dimensional parameter space.
Thus the projective projection center misses the closure of $C$.
The resulting map to $\mathbb P^j$ is finite: its hyperplane
pullback is the ample hyperplane bundle on $C$. Its degree is
$\deg C$.

This map is generically separable. The coefficient functions of $P$
and $z$ generate the function field of $C$, so their differentials
span its $j$-dimensional cotangent space. The evaluation forms span
these linear functions; hence the differentials of $j$ independent
generic evaluation forms are independent. Their generic affine
fiber is therefore reduced and consists of exactly $\deg C$ points.
The same choice works simultaneously for the finitely many actual
$j$-components. Generic target values $\ell_i=c_i$ avoid their
singular loci, intersections with other actual components, and the
separant boundary. The latter is proper after a generic $t_0$ is
chosen, since each actual component has a dense regular locus.
These sections therefore select exactly $J_j$ distinct regular points.

\paragraph{Count the section in initial data.}
Use $t_0$ as Taylor center. The first section is
$u_0=c_0-\lambda_0z$, leaving $d+1$ initial-data variables. The
initial equation has degree at most $q$. The coefficient residuals
have degree at most $K$. Each of the other $j-1$ sections, pulled
back through reconstruction and cleared by $S^\tau$, has degree
at most $T$: its coordinates are coefficient numerators of degree
at most $T$, $zS^\tau$ of degree at most $T$, and $S^\tau$ of
degree at most $T-1$.

These further sections cut the initial hypersurface properly on
$S\ne0$, even before the residuals are imposed. Indeed the Taylor
map there has a linear inverse given by its derivatives through order $d$ and $z$,
so it is an immersion. Evaluation hyperplanes span all its coefficient
coordinates and cut each positive-dimensional image properly for
generic parameters and target values. After the $j-1$ further
sections, the local source dimension at each selected point is at
most $d+1-j$.

Now impose $d+1-j$ generic linear combinations of the residuals.
Their full common zero set already isolates every selected point:
those points lie on no higher-dimensional actual component and the
actual $j$-components have just been cut to points. The pruning
argument of Theorem~\ref{thm:isolated-upper} therefore retains and
isolates all selected points. B\'ezout bounds their number by
$qT^{j-1}K^{d+1-j}$, as claimed. Components on $S=0$ contribute
no selected point and need not be interpreted as actual solutions.

\paragraph{Fixed challenge.}
Work over an algebraic closure of the coefficient field with that
challenge held constant. There are $d+1$ initial variables, initial
equation degree at most $B$, coefficient-numerator degree at most
$b$, and residual degree at most $Bb$. Use the evaluation forms
$P(t_i)$ without a $z$ coordinate. The same projection argument
applies: a nonzero coefficient direction cannot vanish at every
$t$. The first section is $u_0=c_0$. The identical count gives
$B(Bb)^{d-j}b^{j-1}$ for $j\ge1$ and $B(Bb)^d$ for isolated
points. Apply this in particular over the algebraic closure of the
generic challenge field.
\end{proof}

\begin{corollary}[Nonaffine actual curves at every order]
\label{cor:higher-nonaffine-curves}
For a received line on $n$ distinct coordinates and $D<A\le n$,
the actual curve components in Theorem~\ref{thm:higher-actual-components},
other than affine codeword graphs $P=F+zG$, contain at most
\[
 \frac{n-D}{A-D}\,qK^d
\]
pairs with at least $A$ agreements. This is $O(n^d)$ at fixed
jet and challenge degrees, derivative order, and positive gap.
\end{corollary}
\begin{proof}
The agreement-hyperplane count in
Corollary~\ref{cor:nonaffine-actual-curves} uses only actual curve
degrees. Each retained curve has at most $D$ persistent agreement
coordinates, and its nearby-pair count is at most its degree times
$(n-D)/(A-D)$. Sum and apply $J_1\le qK^d$.
\end{proof}

\begin{theorem}[Sharp positive-dimensional and fiber exponents]
\label{thm:positive-component-sharpness}
For every $r\ge1$ and $e\ge1$, put $D=re$ and assume characteristic
zero or $p>D+r-1$. There is an equation of derivative order $r$,
jet degree $r$, challenge degree one, and $X$-degree bounded in terms
of $r$ alone, whose reduced actual regular
solution closure is irreducible of dimension $r+1$ and degree at
least $e^r$. Its reduced challenge fibers have degree $e^{r-1}$.
Thus the exponents $D^d$ and $D^{d-1}$ in
Theorem~\ref{thm:higher-actual-components} cannot be lowered uniformly,
even with a single actual component.
\end{theorem}
\begin{proof}
Set $\phi_j=X^j$ for $1\le j<r$ and $\phi_r=X^r+z$, and define
\[
 Q=\operatorname{Wr}\bigl(eP\phi_1'-P'\phi_1,\ldots,
                           eP\phi_r'-P'\phi_r\bigr).
\]
This is homogeneous of jet degree $r$, and its separant is
\[
 Q_{P^{(r)}}=(-1)^r C_r(eP)^{r-1}
               \bigl(X^r+(-1)^{r-1}z\bigr),
 \qquad C_r=\prod_{j=1}^{r-1}j!.
\]
The proof is the same row calculation as
Equation~\eqref{eq:isolated-separant}. Its coefficient of $P^r$
is $e^r\operatorname{Wr}(\phi_1',\ldots,\phi_r')$, which is
nonzero; hence the jet degree is exactly $r$. The separant proves
exact order $r$, linear highest-derivative degree, and challenge
degree one. Every nonzero polynomial solution is regular.

Each specialized Wronskian column has degree at most $D+r-1<p$.
The polynomial Wronskian criterion from
Theorem~\ref{thm:isolated-lower} therefore says that $Q=0$ precisely
when some nonzero
\[
 H\in\operatorname{span}(X,\ldots,X^{r-1},X^r+z)
\]
satisfies $ePH'-P'H=0$. For nonzero $P$, this makes
$(P/H^e)'=0$. The rational degree of this quotient is at most
$D<p$, so $P=cH^e$ for a nonzero constant $c$. Conversely every
such power satisfies the equation. The zero polynomial also
satisfies it; for $r>1$ it is singular and belongs to the closure
of the nonzero regular locus.

Write $H=\sum_{i=0}^r h_iX^i$ projectively. The condition on $H$
is exactly $h_0=zh_r$. The power map
\[
 \mathbb P^r\longrightarrow\mathbb P^D,\qquad [H]\longmapsto[H^e]
\]
is base-point-free and generically one-to-one, with hyperplane
pullback $\mathcal O(e)$. It is birational onto its image: on the
monic degree-$r$ chart its coefficients are recovered successively
from the coefficients of $H^e$, dividing by $e$, which is invertible.
Its image consequently has degree $e^r$.

The projective subfamily $h_0=zh_r$ is closed over the affine
$z$-line. Its power-map image is closed by properness, and its
relative affine cone is exactly the actual solution closure just
classified. The family of hyperplanes $h_0=zh_r$ is a projective
bundle, so that cone is irreducible of dimension $r+1$. On the dense
monic chart, $z$ is recovered from the constant coefficient of $H$.
Projection forgetting $z$ is thus birational onto the affine cone
of the whole power-map image, of degree $e^r$. Linear projection
cannot increase degree, giving the claimed joint-degree lower bound.
For each fixed $z$, the hyperplane of $H$ is a $\mathbb P^{r-1}$;
its power-map image, and hence its affine cone, has degree $e^{r-1}$.
\end{proof}

\begin{proposition}[Sharp isolated fibers and affine-graph counts]
\label{prop:fixed-fiber-isolated-sharpness}
Fix $d\ge1$ and $D\ge d$, in characteristic zero or $p>D+d$.
There is a challenge-independent equation of order $d$ and jet
degree $d+1$ with exactly
\[
 \binom{D+d+1}{d}=\Theta_d(D^d)
\]
polynomial solutions of degree at most $D$, all regular and isolated.
Adjoining a free challenge turns these into that many actual affine
codeword-graph curve components of degree one. Thus both the $I_0$
and $J_1$ exponents of Theorem~\ref{thm:higher-actual-components}
are sharp as algebraic counts.
\end{proposition}
\begin{proof}
Let $R$ be monic and squarefree of degree $N=D+1$, and take the
Wronskian of
\[
 R(X^j)'-PX^j,\qquad j=0,\ldots,d.
\]
Its separant in $P^{(d)}$ is
$(-1)^{d+1}R^d\prod_{j=1}^d j!$, a nonzero polynomial independent
of $P$. Each column has degree at most $D+d<p$, so the Wronskian
vanishes exactly when $RH'-PH=0$ for some nonzero polynomial
$H$ of degree at most $d$. Hence $P=RH'/H$. For nonconstant $H$,
its roots must belong to $R$; conversely every multiset of at most
$d$ roots gives a solution. Constant $H$ gives $P=0$. The residues
are integers between zero and $d<p$ and distinguish all the multisets.
The number is
\[
 \sum_{s=0}^d\binom{N+s-1}{s}=\binom{N+d}{d}.
\]
Every solution is regular by the separant formula. The finite
classification proves the isolated count; a free challenge gives
one constant-polynomial graph per solution.
\end{proof}

\begin{proposition}[Linear MCA for the logarithmic-derivative family]
\label{prop:logarithmic-family-mca}
For fixed $d$, let $R$ be monic and squarefree of degree $D+1$,
in characteristic zero or $p>d$. Consider the family
\[
 \mathcal P=\{RH'/H:\ H\text{ monic},\ \deg H\le d,
                  \text{ all roots of }H\text{ belong to }R\}.
\]
This includes the zero polynomial. For $D<A\le n$ with
$A-D\ge\eta n$, every received line has $O_{d,\eta}(n)$
full-support bad labels witnessed by $\mathcal P$, for sufficiently
large $n$. A candidate can be reused at many challenges; no prior
assignment of labels is assumed.
\end{proposition}
\begin{proof}
Select one bad candidate at each of $L$ distinct bad labels, and
write $W=P/R$. At each root of $R$ its residue belongs to
$\{0,\ldots,d\}$. On any affine line of graph points $(z,W)$
where $W$ varies, one residue is an injective affine function of
$z$, so the line has at most $d+1$ selected points. If $W$ is
constant, the candidate $P$ is fixed. It has at most $n$ bad labels:
a bad full support must contain an agreeing coordinate with
$g(x)\ne0$, which determines $z$; otherwise $(P,0)$ explains
that support. Thus every graph line has at most
$M=\max\{n,d+1\}$ selected points.

A noncollinear triple gives a nonzero rational relation of numerator
degree at most $C=3d-1$, after clearing its three denominators of
degree at most $d$. It has at most $C$ common agreeing coordinates
outside roots of $R$. There are at most
$(M-2)\binom L2/3$ collinear triples, by counting pairs.
Let $r_0$ count domain roots of $R$, and put $m=n-r_0$ and
$t=A-r_0$. Triple incidence and convexity give
\[
 m\left(\frac{tL}{m}-2\right)_+^3
 \le C L^3+m(M-2)L^2.
\]
For sufficiently large $n$, $t\ge\eta n-1$ and $m\le n$ imply
$t^3>8Cm^2$. If $L<4m/t$ the claim is immediate. Otherwise the
left side is at least $t^3L^3/(8m^2)$, giving
\[
 L\le\frac{8m^3(M-2)}{t^3-8Cm^2}=O_{d,\eta}(n).
\]
This counts genuinely bad labels rather than all nearby challenges,
which may fill the field on a constant candidate graph.
\end{proof}

\paragraph{Verification and MCA scope.}
The exact checker in \path{research/actual_higher_order_components/}
checks the complete solution classification on four small-field fixtures,
the separant through order five, and $40$ reduced sections of the
projected power images and fixed-challenge fibers. These sections
have exactly $e^r$ and $e^{r-1}$ points. It includes a small-
characteristic failure and explicitly checks lower-degree boundary
powers, which must not be lost by using only monic degree-$r$ charts.
The geometric upper bound is a separate written argument. A second
checker exhausts $7{,}203$ polynomials for the fixed-equation family
and verifies its separant through order four.

The positive-dimensional degree saving alone does not prove an
$O(n^d)$ full-MCA bound. Isolated actual joint points can number
$\Theta(D^{d+1})$, and actual affine codeword-graph curve components
can each contribute accidental labels. Their generic challenge
fibers are isolated points, outside the positive-dimensional fiber
bound. The sharpness family here has at most $r$ distinct roots per
candidate and is already covered by the linear-MCA result of
Appendix~\ref{app:bounded-root-mca} in its sufficiently large
characteristic regime. The fixed-gap quadratic MCA lower bound
remains open.
